\chapter{Assessments of the year}
\textbf{Organisation.} Each deadline has two practice papers for training. They announce the expected methods; they may be given at the start of the unit, even if not all the topics have yet been taught. The table indicates the planned time of distribution. Versions A and B are reserved for the assessment.
\ifprof
\paragraph{Marking schemes.} Each exercise is worth 5 points; the values are given question by question. For a 1-point question requiring a proof, award up to 0.5 point for a relevant approach and up to 0.5 point for the justification and conclusion. For a proof by induction marked out of 2 points: property and base case 0.5; hypothesis at a fixed term and inductive step calculation 1; quantified conclusion 0.5. For another value, adapt these proportions. Do not penalise a propagated error twice. The solutions write out the expected reasoning and calculations; accept any correct method within the taught scope.
\paragraph{Duration.} The 2 h tests have four short exercises; the 4 h mock exams have four developed exercises. Adjust only after real feedback from the class. The four 5-point exercises are a choice consistent with the progression and recent papers, without fixing the instructions of the future session.
\fi
\begin{longtable}{@{}p{.17\textwidth}p{.54\textwidth}p{.21\textwidth}@{}}\toprule\textbf{Deadline}&\textbf{Administration guide}&\textbf{Practice papers}\\\midrule\endhead


Test 1&S9 · week of 21 September&From S2\\[3pt]

Test 2&S18 · week of 12 October&From S14\\[3pt]

Common test&S28 · week of 16 November; slot to be harmonised&From S14\\[3pt]

Test 3&S34 · week of 30 November&From S21\\[3pt]

Mock exam 1&S45 · 14 to 16 January, 4 h slot to be fixed&From S36\\[3pt]

Test 4&S57 · week of 8 February&From S44\\[3pt]

Test 5&S68 · week of 22 March&From S52\\[3pt]

Mock exam 2&S71 · 1st to 3 April, 4 h slot to be fixed&From S59\\[3pt]

Final mock exam&S90 · week of 31 May, 4 h slot to be fixed&From S75\\[3pt]

\bottomrule\end{longtable}
